\chapter{Lieux et strates}
\label{chap:lieux}

\section{Les lieux (\emph{locations})}

Un \emph{lieu} est tout endroit physique où vous cultivez : une parcelle entière, une planche dans une parcelle, une serre, un verger, un rang dans un verger. Les lieux sont organisés en \textbf{hiérarchie} : une planche peut être contenue dans une parcelle, elle-même contenue dans une exploitation.

Chaque lieu possède :

\begin{itemize}
  \item Un \textbf{nom} (libre).
  \item Une \textbf{catégorie} (parcelle, planche, serre, verger, rang d'agroforesterie…).
  \item Des \textbf{dimensions} : longueur et largeur en mètres (la surface est calculée automatiquement).
  \item Un éventuel \textbf{parent} pour former la hiérarchie.
  \item Des \textbf{notes} libres.
\end{itemize}

\section{Créer un lieu}

L'écran \textbf{Lieux} présente la hiérarchie des emplacements (les enfants indentés sous leur parent) et un formulaire de création. Pour ajouter un lieu :

\begin{enumerate}
  \item Saisissez le \textbf{nom}.
  \item Choisissez la \textbf{catégorie} (parcelle, planche, serre, verger, rang…). Les catégories marquées « sous abri » (serre, tunnel) comptent dans la courbe d'occupation \emph{sous abri} de l'accueil.
  \item Renseignez la \textbf{longueur} et la \textbf{largeur} en mètres ; la surface est calculée automatiquement.
  \item Sélectionnez éventuellement un \textbf{parent} pour rattacher ce lieu (par exemple une planche dans une parcelle).
  \item Ajoutez des \textbf{notes} si besoin, puis validez.
\end{enumerate}

\section{Les strates}
\label{sec:strates}

Une \emph{strate} décrit une couche de végétation par sa hauteur. Pomone fournit un jeu de strates par défaut (couvre-sol, herbacée basse, herbacée haute, arbustive, arborée haute) que vous pouvez personnaliser.

Chaque strate possède :

\begin{itemize}
  \item Un \textbf{nom} et une \textbf{description}.
  \item Une plage de \textbf{hauteurs min/max} en mètres.
  \item Un \textbf{ordre de tri} pour l'affichage.
\end{itemize}

L'écran \textbf{Strates} liste les strates existantes et un formulaire de création/édition. Chaque ligne porte un bouton \textbf{Supprimer} ; la suppression est refusée (avec un message clair) tant qu'une culture référence la strate — réaffectez ces cultures d'abord.

\section{Vue Carte (occupation des lieux)}

L'occupation des lieux dans le temps fait l'objet d'un écran dédié, la \textbf{Carte} (\emph{Crop Map}) : pour chaque emplacement, une frise des plantations qui l'occupent sur la saison, avec déplacement et division de plantations. Voir le chapitre \ref{chap:carte}.
